% analysis/failure_modes-DE.tex
\section{Fehlermodi und Grenzen der AIM-Pipeline}
\subsection*{Motivation}
Jeder Modellierungs- und Optimierungsrahmen muss neben seinen Fähigkeiten auch
seine Grenzen berücksichtigen. Fehlermodi der AIM-Pipeline entstehen aus dem
Zusammenwirken von Koordinatensystemen, geometrischen Transformationen,
physischer Realisierung und numerischer Optimierung.

\subsection{CRS-bezogene Fehler}
Fehler in Koordinatenreferenzsystemen umfassen falsche CRS-Identifikation,
inkonsistente Transformationen und gemischte Koordinatensysteme in Datensätzen.
\begin{proposition}
CRS-Inkonsistenzen führen zu systematischen räumlichen Verzerrungen, die durch
lokale Alignment-Optimierung nicht korrigiert werden können.
\end{proposition}

\subsection{Projektionsfehler (World--Track)}
Die World--Track-Transformation kann durch mehrdeutige nächste Punkte, Bereiche
hoher Krümmung sowie dünn besetzte oder verrauschte Alignment-Darstellungen
fehlschlagen.
\begin{proposition}
Projektionsfehler führen nichtlokale Verzerrungen der intrinsischen Koordinate
\(s\) ein und beeinflussen alle nachgelagerten Berechnungen.
\end{proposition}

\subsection{Fehlspezifikation des Modells}
Das parametrische Modell kann das tatsächliche Alignment wegen unzureichender
Übergangsfamilien, falscher Parametrisierung oder zu restriktiver Struktur
nicht darstellen. Hybride Korrekturterme mindern, beseitigen dieses Problem
aber nicht.

\subsection{Überanpassung gegenüber Unteranpassung}
Hybride Modelle führen eine Abwägung ein: Bei Unteranpassung ist das
parametrische Modell zu starr; bei Überanpassung erfassen Korrekturterme das
Rauschen.
\begin{proposition}
Regularisierung ist erforderlich, um Modelltreue und Stabilität auszugleichen.
\end{proposition}

\subsection{Realisierungseffekte}
Der Realisierungsoperator führt Krümmungsglättung, Verlust hochfrequenter
Information und Abhängigkeit von Bauprozessen ein.
\begin{proposition}
Bestimmte Merkmale des Ziel-Alignments sind physisch nicht realisierbar und
werden systematisch unterdrückt.
\end{proposition}

\subsection{Instabilität des inversen Problems}
Das inverse Realisierungsproblem ist schlecht gestellt: Mehrere Ziel-Alignments
erzeugen ähnliche realisierte Geometrie, und kleine Datenänderungen können
große Parameteränderungen bewirken.
\begin{proposition}
Regularisierung und Vorwissen sind für stabile Lösungen wesentlich.
\end{proposition}

\subsection{Probleme numerischer Optimierung}
Optimierung kann an schlechten Anfangsschätzungen, nichtkonvexen Zielfunktionen
oder schlecht konditionierten Jacobi-Matrizen scheitern. Trotz strukturierter
dünner Besetzung ist Konvergenz nicht garantiert.

\subsection{Degeneration der Blockstruktur}
Die Blockstruktur kann degenerieren, wenn Kopplungsterme dominieren, Daten die
Variablen unzureichend trennen oder Parameter redundant werden.
\begin{proposition}
Der Verlust der Blockseparierbarkeit verringert Recheneffizienz und
Interpretierbarkeit.
\end{proposition}

\subsection{Messrauschen und Verzerrung}
Messdaten können zufälliges Rauschen, systematische Verzerrung und fehlende
oder inkonsistente Segmente enthalten.
\begin{proposition}
Rauschen pflanzt sich durch die Pipeline fort und kann durch inverse
Operationen verstärkt werden.
\end{proposition}

\subsection{Maßstabsabhängige Fehler}
Auf verschiedenen Maßstäben dominieren unterschiedliche Fehlermodi:
\begin{itemize}
\item großer Maßstab: CRS- und Projektionsfehler,
\item mittlerer Maßstab: Fehlspezifikation des Modells,
\item kleiner Maßstab: Realisierung und Messrauschen.
\end{itemize}

\subsection{Nicht-Eindeutigkeit von Lösungen}
Verschiedene Alignment-Modelle können nahezu identische realisierte Geometrie
erzeugen.
\begin{proposition}
Die Abbildung
\[
\kappa_{\mathrm{target}}\mapsto\gamma_{\mathrm{real}}
\]
ist nicht injektiv.
\end{proposition}

\subsection{Praktische Engineering-Grenzen}
Materialabhängiges Verhalten, langfristige Degradation und Bauvariabilität
liegen außerhalb des Modellumfangs und müssen als externe Bedingungen
behandelt werden.

\subsection{Zentrale Erkenntnis}
\begin{proposition}
Die AIM-Pipeline ist ihrem Wesen nach approximativ und wird sowohl durch
physische als auch durch rechnerische Bedingungen begrenzt.
\end{proposition}

\subsection{Verbindung zu AIM}
\begin{summary}
Fehlermodi der AIM-Pipeline entstehen aus dem Zusammenwirken von Geometrie,
Physik, Daten und Berechnung. Das Verständnis dieser Grenzen ist für robusten
Modellentwurf, stabile Optimierung und zuverlässige Engineering-Anwendung
wesentlich.

Fehleranalyse ist daher kein Hilfsbestandteil, sondern integraler Teil des
AIM-Rahmens.
\end{summary}
